% !TeX root = ../documento.tex
\begin{comment}
Falta Revisar 18/07/2025
\end{comment}
Considering the advances in control strategies and embedded electronics, the enhancement of vehicle safety functions—especially within Advanced Driver Assistance Systems (ADAS)—has increased the demand for predictable and fast braking actuation. Challenges remain regarding interoperability, performance under variable conditions, accuracy, and response time, particularly when the electro-hydraulic unit (ABS/ESC hydraulic control unit) is used as an \textbf{active braking pressure actuator} for longitudinal ADAS functions (e.g., Adaptive Cruise Control, ACC).

In this context, this work proposes and validates, in simulation, a \textbf{digital pressure regulation methodology} for the hydraulic brake module, treating the wheel pressure reference as an input generated by a higher-level controller (e.g., ACC). The validation is conducted in MATLAB/Simulink using a lumped non-linear hydraulic model (two pressures), local linearization for control design, and a discrete-time state-space controller (LQI) with a Luenberger observer, including actuator saturations and a discrete valve logic.

To support future embedded implementation and integration, an electronic prototype platform based on commercial components and the Infineon AURIX TC375 microcontroller was developed. Therefore, simulation evidence and the prototype development indicate the feasibility of the proposed approach and suggest potential application in infrastructures such as ABS, ESC, TCS, ACC, AEB, EBD, and ADAS.

\keywords{Vehicle Brake Control, Hydraulic Pressure Regulation, Digital Control, ADAS, MATLAB/Simulink, TC375}